\section{Conclusion}
\label{sec:conclusion}

We presented ToonSense, a robustness study of multimodal classifiers for Bangla--English memes built on frozen
multilingual text encoders and CLIP. Across CMBan cartoon memes and two larger Bangla meme datasets, we benchmarked
six text encoders, four image encoders and five fusion heads, and subjected the primary XLM-R + CLIP-B/32 models to
twelve linguistic and visual perturbations at three severities.

Four conclusions stand out.
\begin{enumerate}
  \item \textbf{Clean accuracy.} The image modality carries most of the signal, and frozen text encoders are
  outperformed even by character $n$-grams. Fusion adds a small but significant gain once a few thousand training
  memes are available.

  \item \textbf{Romanisation is the main threat.} Writing the same meme in Romanised Bangla is the most damaging
  perturbation on every dataset, costing the fusion models 8--21\% of relative macro-F1. Native-script spelling
  noise and code-mixing are nearly harmless. Fusion halves the damage only when enough training data is
  available, and the perturbed text attracts \emph{more} attribution rather than less.

  \item \textbf{A simple remedy.} Adding transliterated copies of the training memes, with no new labels, reduces the Romanisation drop of
  the fusion models from 7--15\% to 2\% for transliteration and from 8--21\% to 3--5\% for an unseen
  Romanised spelling style, while changing clean macro-F1 by less than one point.

  \item \textbf{Stability is not accuracy.} Visual perturbations flip up to 27\% of predictions with little change
  in macro-F1. Robustness evaluations should therefore report flip rates. CLIP's attention to rendered text is a
  property of cartoon memes rather than of memes in general.
\end{enumerate}

Future work should evaluate fine-tuned and generative vision--language models with the same protocol, collect
naturally Romanised versions of existing memes, verify the contextual/cultural error candidates with native
speakers, and repeat the CMBan analyses once its full label set is released.
